\subsection{Development of the local smoothing problem}\label{subsec:history}

The question grew out of Bourgain's theorem on maximal averages over
circles \cite{Bourgain1986Circles}. At high frequency the Fourier
transform of a circle of radius \(T\) has two leading terms, with phases
\(T|\xi|\) and \(-T|\xi|\). This explains why a spacetime estimate for
waves can help bound the supremum of circular averages over radii.
Sogge used this connection to prove a positive smoothing gain for every
\(p>2\) in two space dimensions \cite[Theorem~4.1]{Sogge1991}. Mockenhaupt, Seeger
and Sogge gave a proof of the circular maximal theorem using angular
square functions \cite{MSS1992}, and developed local smoothing for
Fourier integral operators under a curvature condition on the
associated family of cones \cite{MSS1993}.

A persistent difficulty is to put waves moving in different directions
back together. Wolff decomposed a neighborhood of the cone into thin
angular sectors and controlled their sum by induction on scales. This
gave every gain below \(1/p\) for \(p>74\) in the plane
\cite{Wolff2000}; {\L}aba and Wolff extended the argument to higher
dimensions \cite{LabaWolff2002}. Bourgain and Demeter's cone decoupling
theorem gives this gain for \(p\ge 2(n+1)/(n-1)\) in every dimension
\cite[Theorem~1.2]{BourgainDemeter2015}. A different argument using
radial Fourier multipliers attained the exact gain \(1/p\), with no
additional Sobolev loss, for \(n\ge4\) and \(p>2(n-1)/(n-3)\)
\cite[Theorem~1.2]{HeoNazarovSeeger2011}. In the plane, Guth, Wang
and Zhang reached the critical exponent \(p=4\) by controlling the
\(L^4\) norm of a wave by the spacetime square function of its angular
pieces \cite[Theorems~1.1 and~1.2]{GWZ2020}. For the canonical
sectors of a truncated cone, Cushman, Demeter and Wu give the sharp
power \(1/4\) of the logarithm in this square function inequality
\cite[Theorem~1.4]{CushmanDemeterWu2026}. To obtain a bound in
terms of \(\|f\|_4\), one must also estimate the square function from
the initial data.

In higher dimensions the distinction between a partial gain at \(\pc\)
and a full gain above \(\pc\) is substantial. At \(n=p=3\), Gao, Liu,
Miao and Xi obtained every gain below \(2/9\) by broad Fourier
restriction estimates, compared with the conjectured \(1/3\)
\cite[Corollary~1.3]{GaoLiuMiaoXi2023}. Gan, He, Li and Wu proved the
full gain up to an arbitrarily small loss for \(p\ge10/3\) in
\(\RR^3\) and on compact Riemannian three-manifolds
\cite[Theorem~1.3]{GHLW2026}; see also Wang's study of wave packet
density \cite{Wang2026v4}. Interpolating the \(p=10/3\) estimate with
the \(L^2\) identity gives every gain below \(1/4\) at \(p=3\).
Some three-dimensional metrics require
\(p\ge10/3\) for gains arbitrarily close to \(1/p\)
\cite{MinicozziSogge1997}, whereas the Euclidean exponent is \(3\).
